% 爱德华·康登（综述）
% license CCBYSA3
% type Wiki

本文根据 CC-BY-SA 协议转载翻译自维基百科\href{https://en.wikipedia.org/wiki/Edward_Condon}{相关文章}
\begin{figure}[ht]
\centering
\includegraphics[width=6cm]{./figures/748cfbd4f06fd931.png}
\caption{} \label{fig_ADHkd_1}
\end{figure}
爱德华·尤勒·康登（Edward Uhler Condon，1902年3月2日－1974年3月26日）是一位美国核物理学家、量子力学的先驱之一，并在第二次世界大战期间参与了雷达的研制工作，以及在曼哈顿计划中短暂参与了核武器的开发。以他命名的“弗兰克–康登原理”和“斯莱特–康登规则”均纪念了他在该领域的贡献。\(^\text{[1][2][3]}\)

他曾于1945年至1951年间担任美国国家标准局（现为美国国家标准与技术研究院，NIST）的第四任局长。1946年，康登出任美国物理学会会长；1953年，他担任美国科学促进会会长。

在麦卡锡主义时期，康登是最早成为美国众议院非美活动调查委员会攻击目标的知名科学家之一。1948年，他被公开指控为“美国原子安全体系中最薄弱的一环”，理由是他掌握大量机密信息、与原子弹研制有关联，并被指涉嫌同情共产主义和苏联。康登事件成为当时反对麦卡锡主义者（尤其是科学界）关注的标志性案件之一，舆论广泛关注，他也得到了包括美国总统哈里·杜鲁门在内的众多知名科学家的公开支持。\(^\text{[4]}\)

1968年，康登作为《康登报告》的主要作者再次广为人知。该报告是由美国空军资助的官方评估，得出的结论是“不明飞行物都有平凡的解释”。月球上的“康登环形山”即以他的名字命名，以纪念他的贡献。
\subsection{背景}
\begin{figure}[ht]
\centering
\includegraphics[width=6cm]{./figures/3223424e58a1eb6d.png}
\caption{图1. 弗兰克–康登原理能量图。由于电子跃迁相比于原子核运动要快得多，因此在跃迁过程中，最有利的振动能级是那些对应于原子核坐标变化最小的能级。图中所示的势阱表明，跃迁在 $v = 0$与$v = 2$之间最为有利。} \label{fig_ADHkd_2}
\end{figure}
爱德华·尤勒·康登于1902年3月2日出生在新墨西哥州的阿拉莫戈多，父母分别是威廉·爱德华·康登和卡罗琳·尤勒。他的父亲当时正在监督一条窄轨铁路的建设，\(^\text{[5][6]}\)当地的许多铁路都是由伐木公司修建的。1918年，康登在加利福尼亚州奥克兰完成高中学业后，曾在《奥克兰询问报》及其他报纸担任记者三年。\(^\text{[5]}\)他具有爱尔兰血统。\(^\text{[7]}\)

随后，他进入加利福尼亚大学伯克利分校就读，最初加入化学院；后来得知他高中时的物理老师成为了该校教师后，他转而改修理论物理。\(^\text{[8]}\)康登仅用三年时间获得学士学位，又用两年取得博士学位。\(^\text{[5]}\)他的博士论文结合了雷蒙德·塞耶·伯奇在带状光谱强度测量与分析方面的研究，以及詹姆斯·弗兰克的一个建议。\(^\text{[9][10]}\)

凭借国家研究委员会奖学金，康登前往德国哥廷根大学师从马克斯·玻恩，又赴慕尼黑大学师从阿诺德·索末菲。在后者指导下，康登以量子力学的形式重写了自己的博士论文，并提出了著名的“弗兰克–康登原理”。\(^\text{[9]}\)

1927年秋，他在《物理评论》上看到一则招聘广告后，加入贝尔电话实验室从事公共关系工作，主要负责宣传他们关于电子衍射的发现。\(^\text{[5][11]}\)在贝尔实验室工作期间，他还研究了词语使用频率的分布，并将结果与心理学中的韦伯–费希纳定律联系起来。\(^\text{[12]}\)
\subsection{职业生涯}
\subsubsection{早期经历}
\begin{figure}[ht]
\centering
\includegraphics[width=6cm]{./figures/6f359ea6e05a3107.png}
\caption{阿诺德·索末菲，德国理论物理学家，曾在20世纪20年代指导过康登。} \label{fig_ADHkd_3}
\end{figure}
康登曾在哥伦比亚大学短暂任教，并于1928年至1937年期间担任普林斯顿大学物理学副教授，\(^\text{[5]}\)其间曾在明尼苏达大学任教一年。\(^\text{[13]}\)他与菲利普·M·莫尔斯合著了《量子力学》，这是1929年出版的第一本英文版量子力学教材。随后，他与G.H.肖特利合著了《原子光谱理论》，该书自1935年出版起便被誉为该领域的“权威圣经”。\(^\text{[14][15][16][Note 1]}\)

1937年起，康登担任位于匹兹堡的西屋电气公司研究副主任，在那里他创建了核物理、固体物理和质谱学的研究项目。随后，他领导公司在微波雷达开发方面的研究工作。\(^\text{[13]}\)他还参与了用于分离原子弹原料铀的设备研发。\(^\text{[5]}\)在第二次世界大战期间，康登担任国家防务研究委员会的顾问，并协助组织了麻省理工学院的辐射实验室。\(^\text{[13][15]}\)1940年5月11日，康登在纽约世界博览会上展示了名为“Nimatron”的机电装置，这是一台用于玩尼姆游戏的早期计算机式设备。\(^\text{[18]}\)
\subsubsection{政府任职}
\begin{figure}[ht]
\centering
\includegraphics[width=6cm]{./figures/1ec8df323ef9ca2e.png}
\caption{J·罗伯特·奥本海默约摄于1944年在第二次世界大战期间领导了“曼哈顿计划”，而康登曾在20世纪40年代初短暂参与该计划的工作。} \label{fig_ADHkd_4}
\end{figure}
1943年，康登加入了“曼哈顿计划”。然而，仅仅六周后，他便因与项目军事负责人莱斯利·R·格罗夫斯将军在安全问题上的冲突而辞职。格罗夫斯反对康登的上级J·罗伯特·奥本海默与芝加哥大学冶金实验室主任进行讨论。[19]在辞职信中，康登表示他认为留在西屋电气公司对战争贡献更大，并解释说，洛斯阿拉莫斯实验室的安全措施令他极为不安，甚至可能使他陷入抑郁、难以有效工作。[20] 康登对奥本海默没有反抗格罗夫斯感到失望，却不知道当时奥本海默自己也尚未获得安全许可。[19]

自1943年8月至1945年2月，康登在加州大学伯克利分校兼职担任顾问，从事铀-235与铀-238的分离研究。[21] 1944年，康登当选为美国国家科学院院士。[16]1945年6月，康登受邀与多位美国著名科学家一同前往莫斯科，出席俄罗斯科学院成立220周年庆典。他表达了强烈的出席意愿。然而，当格罗夫斯得知此事后，立即联系西屋电气公司，称他认为康登此行可能存在泄露仍在进行的原子弹研究机密的风险。对此，康登试图通过白宫提出抗议。随后，格罗夫斯要求美国国务院撤销康登的护照，该请求最终被执行。[22]

战后，康登在组织科学家推动原子能由民用机构而非军事机构管理方面发挥了重要作用，[23] 主张以开放而非严格保密的方式推动原子能的发展。他担任参议员布赖恩·麦克马洪的科学顾问——麦克马洪是参议院原子能特别委员会主席，该委员会起草了《麦克马洪–道格拉斯法案》。该法案于1946年8月通过，设立了美国原子能委员会，将原子能的管理权交由民用机构。[5][15][23]康登秉持国际主义立场，主张国际科学合作，并加入了美苏科学协会。[24] 1947年，他当选为美国艺术与科学学院院士。[25]

时任美国商务部长、前副总统亨利·A·华莱士在1945年10月结识康登后，向总统推荐他出任美国国家标准局（National Bureau of Standards，现为美国国家标准与技术研究院 NIST）局长。总统哈里·S·杜鲁门接受了这一建议并提名他担任该职，参议院也一致通过了任命。康登自1945年起担任标准局局长，直至1951年离任。[26][5][21]此外，康登在1946年担任美国物理学会会长，[15][16] 并与独立公民艺术、科学与职业委员会保持联系或为其成员之一。[27] 1949年，他当选为美国哲学学会院士。[28]
\subsubsection{抨击}
\textbf{20世纪40年代}
\begin{figure}[ht]
\centering
\includegraphics[width=6cm]{./figures/da962a6675a86ce0.png}
\caption{} \label{fig_ADHkd_5}
\end{figure}
在20世纪40年代，康登的安全许可多次遭到质疑、审查，并最终被重新确认。

1946年5月29日，联邦调查局（FBI）局长约翰·埃德加·胡佛致信总统杜鲁门，信中点名多位高级政府官员涉嫌参与苏联间谍网络，并称康登“完全就是一个伪装的间谍”。（数十年后，参议员丹尼尔·帕特里克·莫伊尼汉称此指控为“毫无根据的走廊流言”。）杜鲁门政府无视了胡佛的这些指控。[29]

1947年3月21日，杜鲁门签署了美国总统行政命令第9835号，即“忠诚令”（Executive Order 9835，又称 Loyalty Order）。[30] 同月，众议员J·帕内尔·托马斯，即众议院非美活动调查委员会主席，向《华盛顿时报先驱报》提供了不利于康登忠诚度的信息，导致该报发表了两篇质疑他的报道。[31][32]托马斯选择将康登作为重点攻击目标有多重原因：他本就对科学界的国际主义精神缺乏同情，同时希望借此争议推动增加自己委员会的预算，削弱康登所支持的《麦克马洪法案》，并在选举季博取有利的媒体关注。[33]最终，美国商务部于1948年2月24日正式为康登洗清了“不忠诚”的指控。[34]
\begin{figure}[ht]
\centering
\includegraphics[width=6cm]{./figures/7e461019ca8c37fe.png}
\caption{J·帕内尔·托马斯（J. Parnell Thomas，摄于1939年）在1947年和1948年期间攻击了康登的声誉。} \label{fig_ADHkd_6}
\end{figure}
尽管如此，1948年3月2日发布的一份众议院非美活动调查委员会报告仍声称：“康登博士似乎是我国原子安全体系中最薄弱的一环。”[32][35]对此，康登回应道：“如果我真的是原子安全体系中最薄弱的一环，那倒是令人欣慰的事——因为这意味着国家可以完全放心。我是一个绝对可靠、忠诚、认真并全心全意为国家利益服务的人，这一点，我的整个生命与职业生涯都已清楚地证明。”[36]1948年3月3日，新墨西哥州民主党参议员丹尼斯·查韦斯在《国会记录》中宣读了记者马奎斯·查尔兹的一篇文章。文章指出，对康登的不利行动基于极其站不住脚的证据，而且很可能出于政治动机——尤其是与前商务部长、提名康登担任国家标准局局长的亨利·华莱士之间的政治分歧有关。[37]

1948年3月5日，来自明尼苏达州的共和党众议员乔治·麦金农在国会发言中表示：“议长先生，今日的报纸报道称，商务部长W·艾弗雷尔·哈里曼拒绝回应国会的传票，拒绝提供有关康登博士的信息。我无意对该事件的具体事实作出判断，但我必须指出，这种做法与本届政府在整个会期内对待国会传票所采取的保密态度如出一辙。”[35]1948年3月6日，《华盛顿邮报》发表社论称：“商务部长拒绝移交其部门忠诚委员会关于爱德华·U·康登博士的档案，是有充分先例可循的。”报社还反对另一项提议——将康登案的档案提交给公务员委员会下属的“忠诚审查委员会”。《邮报》指出，商务部自己的忠诚委员会已经为康登洗清了指控，该决定理应维持不变。[38]1948年3月8日，加利福尼亚州民主党众议员切斯特·E·霍利菲尔德指出：“……我想提醒各位议员注意我于1947年12月4日提交的H.R. 4641号法案。该法案旨在规范国会调查委员会的调查程序，并保障被调查人员的合法权利。如果该法案得以通过，那么像E.U.康登博士这样享誉世界的科学家，也能获得与偷鸡贼或违反交通法规者同样的权利——也就是有权为自己辩护。打着国会豁免权幌子的‘人格暗杀’，无论来自国会议员还是国会委员会，都是一种危险而可憎的嘲弄。”[35]1948年3月9日，爱达荷州民主党众议员、当时进步党副总统候选人格伦·H·泰勒，纽约州民主党众议员伊曼纽尔·塞勒，以及纽约州美国劳工党众议员利奥·伊萨克森均发表声明，谴责众议院非美活动调查委员会对康登持续不断的攻击。[35]
\begin{figure}[ht]
\centering
\includegraphics[width=6cm]{./figures/3b1ef82a1caedae1.png}
\caption{阿尔伯特·爱因斯坦（Albert Einstein，摄于1947年）为康登进行了辩护。} \label{fig_ADHkd_7}
\end{figure}
康登的支持者中包括阿尔伯特·爱因斯坦和哈罗德·尤里。哈佛大学整个物理系以及众多专业组织都曾致信杜鲁门总统，为康登辩护。[39]1948年4月12日，“原子科学家紧急委员会”（举办了一场晚宴，以表明对康登的支持，活动的发起人中共有九位诺贝尔奖得主。[40]相比之下，美国国家科学院仅考虑发表一份批评众议院非美活动调查委员会调查程序的声明，而非直接为康登辩护。尽管在院士内部该声明获得了压倒性支持（275票赞成，35票反对），但科学院领导层最终决定不公开发布声明，而是选择私下与托马斯议员沟通。[41]1948年7月15日，美国原子能委员会批准了康登的安全许可，使他能够在国家标准与技术研究院接触机密信息。[42]

1948年9月，在美国科学促进会年会上，总统杜鲁门发表讲话，康登就坐于主席台附近。杜鲁门在演讲中严厉谴责托马斯议员及众议院非美活动调查委员会（HUAC），指出这种做法可能导致“在一种人人都担心被无端传言、流言蜚语和诽谤公开曝光的氛围中，重要的科学研究将无法继续进行”。他称HUAC的活动是“我们当今所面临的最不美国化的事情——这种氛围更像是一个极权国家”。[5]康登坚决反对国会试图在科学界内部进行所谓“安全风险”排查的做法。1949年6月，他在写给奥本海默的一封措辞严厉的信中批评后者向HUAC提供了有关同事的信息，信中写道：“我为此夜不能寐，怎么也想不通你怎能这样谈论一个你认识多年的朋友，一个你非常清楚他既是优秀物理学家又是好公民的人。”[19]1949年7月，康登在参议院一个就审议委员会运作规则的分委员会上作证。他批评托马斯和HUAC举行秘密听证会，并随后泄露信息以诋毁他和其他科学家的忠诚度。他指出，该委员会拒绝了他本人及其同事要求公开听证、以便自我辩护的请求。[43]

1949年，著名记者爱德华·R·默罗邀请同事唐·霍伦贝克参与电台旗舰频道WCBS播出的创新性媒体评论节目《CBS评析新闻界》。在节目中，霍伦贝克讨论了爱德华·U·康登、阿尔杰·希斯和保罗·罗伯逊等人物。[44]在谈到康登时，霍伦贝克批评了一些反共人士的做法，认为他们的方式完全错误：

“共产主义者最希望的，莫过于被与那些热爱自由的非共产主义者（如康登）混为一谈……这样一来，他们就更容易掩饰自己的真实本性，并声称‘共产主义者’这个词毫无意义……但与此同时，我们也不能因为有人滥用揭露之名，就因此停止对真正的共产主义者进行正当的揭露。”[44]

\textbf{20世纪50年代}
\begin{figure}[ht]
\centering
\includegraphics[width=6cm]{./figures/3667d3cef2ebe2bb.png}
\caption{哈里·S·杜鲁门总统（Harry S. Truman，图中正在签署宣布国家进入紧急状态并授权美国参加朝鲜战争的宣言）于1945年提名康登担任美国国家标准局局长。} \label{fig_ADHkd_8}
\end{figure}
1951年，在历经多年调查与审查后，康登的名誉终于得以彻底澄清。同年，他离开政府部门，出任康宁玻璃工厂（Corning Glass Works）研发部门负责人。他表示，自己之所以离职，是因为政府职位的年薪仅为14,000美元。杜鲁门总统发表声明，对康登的贡献予以高度赞扬：“您在极为关键的岗位上，以始终如一的忠诚与专注履行了作为局长的职责。凭借您在科学界的崇高声望，以及对局内事务的卓越管理，您使这一机构变得前所未有的重要。”两位共和党众议员声称，康登因被调查为“安全风险”而“在压力下离职”，对此商务部长查尔斯·索耶予以否认。[45]

1951年，康登担任华盛顿哲学学会会长。[46] 1951年12月27日，他当选为1953年度美国科学促进会会长。[47][48][Note 2]1952年9月，康登在国会委员会作证时，首次在宣誓条件下正式否认了此前针对他所有“不忠诚”指控。[31] 然而，众议院非美活动调查委员会在其1952年度报告中仍声称，康登不适合获得安全许可，理由是他“倾向于与不忠或忠诚可疑的人交往，并对必要的安全规定抱有轻视”。[49]1952年12月30日，康登在美国科学促进会年会上正式就任会长。据《原子科学家公报》报道：“在他宣誓就职时，现场的热烈掌声体现了会员们对他忠诚与正直的再次信任与肯定。”[48]

五个月后，按照惯例，康登在离开政府部门时，其安全许可被撤销。[31][48] 1954年7月12日，他再次获得安全许可，该消息于10月19日公布，但仅两天后的10月21日，美国海军部长查尔斯·S·托马斯便暂停了该许可。[31]副总统理查德·尼克松声称自己为此负责，而芝加哥原子科学家组织则指责此举是“对国家安全制度的政治滥用”；不过托马斯部长否认尼克松参与其中。[50][51]随后，康登撤回了安全许可申请，并于12月辞去在康宁公司的职位，因为该公司正寻求政府科研合同，而他由于缺乏安全许可，无法参与军事研究。回顾多年经历的安全审查后，他表示：“我不愿再继续面对这一无休止的、反复的审查过程。”[31] 在康宁公司任职期间，康登的安全许可相关法律费用均由公司承担。[52]

1958年，康登回忆自己当年的决定时写道，他之所以辞职，是因为他认为艾森豪威尔政府“在政策上纵容、甚至默许了对科学家的迫害，或至少对他人攻击和诋毁科学家的行为表现出冷漠无情的态度。我认为形势已无可挽回，而我已经在长达七年的抗争中，尽了一个人所能尽的一切努力。”[53]
\begin{figure}[ht]
\centering
\includegraphics[width=6cm]{./figures/ea9f8f299c940d78.png}
\caption{卡尔·萨根曾讲述过康登与“忠诚审查委员会”之间的一次交锋。} \label{fig_ADHkd_9}
\end{figure}
多年后，卡尔·萨根曾讲述康登描述的一次“忠诚审查委员会”听证经历。委员会的一位成员表示担忧道：“康登博士，这里写道，您曾站在物理学中一个名为……量子力学的革命性运动的最前沿。本次听证会认为，如果您能够领导一场科学革命，那么您也可能领导另一场革命。”康登回答说：“我相信阿基米德在公元前三世纪提出的‘浮力原理’；我相信开普勒在十七世纪发现的‘行星运动定律’；我相信牛顿的运动定律）……”随后，他继续列举了历代科学家的成果，包括伯努利、傅里叶、安培、玻尔兹曼和麦克斯韦。[54]他私下曾说过一句颇具讽刺意味的话：“我加入任何目标看起来高尚的组织，我不会去问里面有没有共产党员。”[55]
\subsubsection{学术生涯}
康登于1956年至1963年间担任圣路易斯华盛顿大学物理学教授，之后自1963年起在科罗拉多大学博尔德分校任教，并兼任联合实验天体物理研究所研究员，直至1970年退休。[5]

1966年至1968年间，康登主持了博尔德大学的不明飞行物研究项目，即著名的“康登委员会”。他之所以被选中，是因为他在学界的崇高地位以及他此前从未公开表态过自己对UFO问题的立场。康登后来写道，他同意领导该项目，是“基于美国空军科学研究办公室的请求，出于履行社会公共职责的考虑”。[56] 该委员会于1968年发布最终报告——报告综合了美国空军“蓝皮书计划”的资料以及两个民间组织收集的报告——得出的结论是：不明飞行物都有平凡的、自然的解释。[57]

康登还曾担任美国物理学学会和美国物理教师协会主席（1964年）。[15] 他于1968年至1969年间出任科学社会责任协会主席，并于1970年担任全国理性核政策委员会共同主席。[15]他与亚利桑那大学的休·奥迪肖共同编辑了《物理学手册》。[5] 1968年，他获得了由光学学会颁发的弗雷德里克·艾夫斯奖章。[58]在他退休时，同事们以出版论文纪念集的方式向他致敬。[59]
\subsection{全球政策}
康登是发起召开“世界宪法制定大会”协议的签署者之一。[60][61]由此，人类历史上首次召开了旨在起草并通过《地球联邦宪法》的世界制宪大会。[62]
\subsection{个人生活与逝世}
1922年，康登与出生于芝加哥的艾米莉（或艾玛）·洪齐克（Emilie/Emma Honzik，1899–1974）结婚。她是一名捷克语翻译。他们育有两子一女。[63] 他们的儿子约瑟夫·亨利·康登（Joseph Henry Condon，1935年2月15日－2012年1月2日）是一位物理学家（获西北大学博士学位）兼工程师，曾在贝尔实验室从事数字电话交换机的研究，并共同发明了贝尔国际象棋计算机。[63]

康登是一名贵格会教徒，并自称为“自由主义者”。[26]

1974年3月26日，康登在科罗拉多州博尔德社区医院去世，距他72岁生日仅24天。[5] 他的妻子艾玛在7个多月后去世。[64]
\subsection{遗产与影响}
\begin{figure}[ht]
\centering
\includegraphics[width=6cm]{./figures/644172e41bda2779.png}
\caption{由“月球轨道探测器1号”拍摄的康登环形山（NASA/L&PI 图像）} \label{fig_ADHkd_10}
\end{figure}
美国国家标准与技术研究院每年都会颁发以康登命名的“康登奖”。该奖项旨在表彰在NIST内部科学与技术书面阐述方面取得卓越成就的个人。该奖自1974年设立以来，已成为该机构的重要荣誉之一。[65]

月球上的“康登环形山”亦以他的名字命名，以纪念他对科学的贡献。[66]

康登的同事刘易斯·布兰斯科姆曾评价道：“‘水门事件’对爱德华·康登而言并不意外，其后续影响也在他意料之中。我想，他一定希望能活着看到国会弹劾调查的结果。但康登深知，自由是要付出代价的，而他也已为此贡献了自己的一份力量。他的理想主义、幽默感与无穷的活力，使他那对更美好世界的不懈追求，显得像是一种乐观的信念。”[67][68]

在克里斯托弗·诺兰执导的2023年电影《奥本海默》中，演员奥利·哈斯基维饰演了爱德华·康登。[69]
\subsection{主要著作}
\begin{itemize}
\item Condon, E. U.; Shortley, G. H.（1935）. The Theory of Atomic Spectra（《原子光谱理论》）. 剑桥大学出版社。[70]
\item Condon, E. U.; Odabaşı, Halis（1980）. Atomic Structure（《原子结构》）. 剑桥大学出版社。[71][72]
\end{itemize}
\subsection{另见}
\begin{itemize}
\item 罗纳德·威尔弗雷德·格尼
\item 量子隧穿
\item 麦卡锡主义
\item 反共主义
\end{itemize}
\subsection{注释}
\begin{enumerate}
\item 一些资料将《原子光谱理论》的出版年份记为1936年，但影印本版本确认1935年才是正确的版权年份，与惠勒的记录一致。[17]
\item 《纽约时报》称康登应于1954年出任该组织主席，但王在《Security》一文第265页中确认其任期实际为1953年。[47][48]
\end{enumerate}
\subsection{参考文献}
\begin{enumerate}
\item Edward Condon（已于2013年11月5日存档于 Wayback Machine）于1944年当选为美国国家科学院院士。
\item “APS Fellow Archive.” [www.aps.org](http://www.aps.org). 于2016年3月3日存档。检索于2018年4月22日。
\item Branscomb, Lewis M.（1974年6月），“Edward Uhler Condon”，Physics Today，第27卷第6期，页68–70。Bibcode: 1974PhT....27f..68B。doi:10.1063/1.3128661。
\item Wang 1992, 第238、249、256脚注47页。
\item M’Ehle, Victor K.（1974年3月27日），“Edward Condon, Leader In A-Bomb Creation, Dies”，The New York Times. 于2012年11月6日存档。检索于2012年10月16日。
\item Morse 1976，第125页。
\item 联邦调查局. Edward Uhler Condon.
\item Morse 1976，第126页。
\item Morse 1976，第127页。
\item Condon, Edward Uhler（1927）. On the theory of intensity distribution in band systems（博士论文）. 加利福尼亚大学伯克利分校。OCLC 21068286. ProQuest 301737248. 检索于2019年5月15日。
\item Morse 1976，第128页。
\item Edward Condon（1928年3月1日），“Statistics of Vocabulary”，Science，第67卷第1733期，页300。doi:10.1126/SCIENCE.67.1733.300。ISSN 0036-8075。PMID 17782935。Wikidata Q81031549。
\item Wang 1992，第242页。
\item Wheeler 1998，第122页。
\item Branscomb, Lewis M. “Edward U. Condon, 1902–1974.” Washington University Library. 于2013年1月5日存档。检索于2012年10月16日。
\item Wang 1992，第241页。
\item Condon, E. U.; Shortley, G. H.（1935）. The Theory of Atomic Spectra. 剑桥大学出版社。ISBN 978-0521092098。于2014年6月27日存档。检索于2016年3月31日。
\item “1940: Nimatron.” platinumpiotr.blogspot.com. 于2018年12月25日存档。检索于2018年4月22日。
\item Bird & Sherwin 2005，第223–224页。
\item Kelly, Cynthia C.（编）（2007）. The Manhattan Project: The Birth of the Atomic Bomb in the Words of its Creators, Eyewitnesses, and Historians.原子遗产基金会。页137–138。
\item Wang 1992，第243页。
\item Wang 1992，第262–263页。
\item Wang 1992，第243–234页。
\item Wang 1992，第244页、244脚注15。
\item “Edward Uhler Condon.” American Academy of Arts & Sciences. 2023年2月9日。检索于2023年3月2日。
\item Wang 2001。
\item “Edward U. Condon Papers.” American Philosophical Society. 1998年。检索于2019年10月20日。
\item “APS Member History.” search.amphilsoc.org. 检索于2023年3月2日。
\item Moynihan, Secrecy, 页63–68。
\item Harry S. Truman, Executive Orders.（已于2017年12月31日存档于 Wayback Machine）The Federal Register, U.S. National Archives。
\item “康登放弃安全许可之争”。《纽约时报》，1954年12月14日。于2018年7月22日存档。检索于2012年10月16日。
\item Wang 1992，第246页。
\item Wang 1992，第252–255页。
\item 《第80届国会第二次会议的议程与辩论》。《国会记录》，第94卷。美国政府印刷局，1948年，第4779页。
\item 《国会记录》。美国政府印刷局，1948年3月，第1987页（兰金报告），第2018页（查韦斯声明），第2222页（麦金农声明），第2337页（H.R. 4641 / 霍利菲尔德发言），第2389页（泰勒），第2405页（塞勒、兰金），第2407页（伊萨克森），第2407–2408页（霍利菲尔德）。检索于2019年10月20日。
\item Wang 1992，第248–249页。
\item Childs, Marquis（1948年3月3日），“华盛顿来电：针对康登的抹黑”。《华盛顿邮报》，第2018页。检索于2019年10月20日。
\item “忠诚档案”。《华盛顿邮报》，1948年3月6日，第8页。
\item Wang 1992，第249页。


\item Wang 1992，第249–250页。
\item Wang 1992，第251页。
\item Wang 1992，第255页。
\item Knowles, Clayton（1949年8月21日），“康登抨击国会调查中的‘泄密’”（*Condon Hits 'Leaks' in House Inquiries*）。《纽约时报》（*The New York Times*）。于2018年7月22日存档。检索于2012年10月16日。
\item Ghiglione, Loren（2008）.*CBS的唐·霍伦贝克：麦卡锡主义时代的诚实记者*（*CBS's Don Hollenbeck: An Honest Reporter in the Age of McCarthyism*）。哥伦比亚大学出版社（Columbia University Press），第146–147页（康登），147页（反击），148–149页（希斯），149–150页（罗伯逊）。ISBN 978-0231516891。检索于2015年9月10日。另见“希斯–钱伯斯案”（*Hiss Chambers*）。
\item “康登博士为更高薪水辞职”。《纽约时报》，1951年8月11日。于2012年11月6日存档。检索于2012年10月16日。
\item Condon, Edward（1960）.“六十年的量子物理”。华盛顿哲学学会公报，第16卷，第83页。
\item “康登博士当选科学家组织主席”。《纽约时报》，1951年12月28日。于2017年3月16日存档。检索于2012年10月16日。
\item Wang 1992，第265页。
\item Wang 1992，第264–265页。
\item “尼克松警告党内红色分子之敌”。《纽约时报》，1954年10月23日。于2018年7月22日存档。检索于2012年10月16日。
\item “尼克松在康登事件中的言论被引用”。《纽约时报》，1954年12月17日。于2018年7月22日存档。检索于2012年10月16日。
\item Bird & Sherwin 2005，第460页。
\item Wang 1992，第266页。
\item Sagan 1996，第248–249页。
\item Wheeler 1998，第113页。
\item Dick 1996，第293页。
\item “用大锤敲坚果”。《自然》，第221卷，1969年3月8日，第899–900页。doi:10.1038/221899a0。
\item “弗雷德里克·艾夫斯奖章 / 奎因奖”。光学学会。于2014年2月1日存档。检索于2012年10月16日。
\item Brittin, Wesley E.; Odabasi, Halis（编）（1971）.现代物理专题：向爱德华·U·康登致敬。伦敦：Hilger出版社。
\item “塞恩·里德致信海伦·凯勒，请求其签署《世界宪章》以支持世界和平，1961年”。海伦·凯勒档案馆，美国盲人基金会。检索于2023年7月1日。
\item “《世界宪法协调委员会致海伦的信，附当前资料》”。海伦·凯勒档案馆，美国盲人基金会。检索于2023年7月3日。
\item “制定地球宪法 | 全球战略与解决方案”。《世界问题百科全书》 国际协会联盟。于2023年7月19日存档。检索于2023年7月15日。
\item “贝尔国际象棋计算机的发明者逝世，享年76岁”。NJ.com，2012年3月2日。于2014年12月17日存档。检索于2014年11月17日。
\item “康登夫人讣告”。《纽约时报》，1974年11月4日，第40页。
\item “物理实验室：2005–2007年奖项与荣誉”。美国国家标准与技术研究院。于2014年2月2日存档。检索于2012年10月16日。
\item “康登环形山”。《行星命名地名录》，国际天文学联合会。于2012年7月13日存档。检索于2012年10月16日。
\item Morse 1976，第143页。
\item Branscomb, Lewis M.（1974）. “Edward Uhler Condon.” Physics Today, 第27卷第6期，页68–70。doi:10.1063/1.3128661。
\item Tinoco, Armando（2023年8月29日）. “克里斯托弗·诺兰将《奥本海默》拍摄周期压缩至57天，以重现洛斯阿拉莫斯”。Deadline. 检索于2025年5月31日。
\item Judd, B. R.（1975）. “〈原子光谱理论〉的若干观点”。载于 G. Putlitz、E. W. Weber、A. Winnacker（编），《原子物理学 4》，普雷纳姆出版社，页13–17。doi:10.1007/978-1-4684-2964-0_3。ISBN 978-1-4684-2966-4。
\item “Halis Odabaşı 传记”。于2020年2月20日存档。检索于2020年2月2日。
\item Collins, P. D. B.（1981）. “书评：《原子结构》”。Physics Bulletin，第32卷，第80页。doi:10.1088/0031-9112/32/3/035。

\end{enumerate}